% !TEX program = xelatex
% ============================================================
% 推导文档（工作文件，不进论文正文）
% 主题：直方图融合系数 alpha ∝ 1/(N̂+N̂₀) 的完整推导与数值验证
% 阅读路线：§1 结论 → §4 桥 → §7 系数；细节需要时再进 §5–§6、§8
% ============================================================
\documentclass[11pt]{ctexart}
\usepackage[a4paper,margin=2.4cm]{geometry}
\usepackage{amsmath,amssymb,booktabs}
\newcommand{\E}{\mathbb{E}}
\newcommand{\Var}{\operatorname{Var}}

\title{直方图融合系数 $\alpha\propto 1/(\hat N+\hat N_0)$：\\完整推导与数值验证}
\date{2026 年 10 月 23 日（工作文档，不进论文正文）}

\begin{document}
\maketitle

\section{问题与结论}

融合管线是：邻居 $q$ 的直方图经对齐（平移 $\hat k_{pq}$）得到 $\widetilde F_{pq}(b)$，乘上系数 $\alpha_{pq}$ 后相加得到融合直方图，再从融合直方图提取回波中心作为深度估计。本文档只回答一个问题：\textbf{$\alpha_{pq}$ 应当取什么形式？}

结论先行。记 $\hat N_q$ 为邻居 $q$ 的信号光子数（直方图峰显著度）、$\sigma_h$ 为 IRF 宽度、$\sigma_\epsilon$ 为同表面对齐残差的标准差，则
\begin{equation}
  \alpha_q\;\propto\;\frac{1}{\hat N_q+\hat N_0},
  \qquad
  \hat N_0\;\triangleq\;\frac{\sigma_h^2}{\sigma_\epsilon^2}.
  \label{eq:conclusion}
\end{equation}

配套两句解读（理解全文的关键）：
\begin{itemize}
  \item \textbf{有效位置权重} $w_q\triangleq\alpha_q\hat N_q=\dfrac{\hat N_q}{\hat N_q+\hat N_0}$：随光子数\textbf{递增、饱和}于 $1$。饱和的来源是"位置精度被对齐残差 $\sigma_\epsilon$ 封顶"——光子再多，位置也不会比对齐残差更准。
  \item \textbf{直方图系数} $\alpha_q$：\textbf{递减}。它不是另一个物理因子，就是有效位置权重\textbf{除以幅度}之后的形态。
\end{itemize}

换言之：精度因子从来没有被"幅度因子替代"。$w_q=\alpha_q\hat N_q$ 是定义式的换算——两个坐标描述同一个权重，一个在"位置"坐标里（递增饱和），一个在"直方图"坐标里（递减）。

\section{记号（本文档局部）}

\begin{center}
\begin{tabular}{@{}llp{0.46\textwidth}@{}}
\toprule
符号 & 含义 & 备注 \\
\midrule
$b$ & bin 坐标 & \\
$h(\cdot)$ & 归一化 IRF，宽度 $\sigma_h$ & 取高斯便于展开；一般形状见 \S9 \\
$\hat N_q$ & 邻居 $q$ 的信号光子数 & 鼓包幅度；实测用峰显著度估计 \\
$\mu_p$ & 待估的公共回波中心 & \\
$\epsilon_q$ & 邻居 $q$ 的对齐残差 & $\mu_q=\mu_p+\epsilon_q$，独立，$\Var=\sigma_\epsilon^2$ \\
$v_q$ & 仅计数噪声下 $q$ 的中心估计方差 & $v_q=\sigma_h^2/\hat N_q$（\S5.2） \\
$v^{\mathrm{eff}}_q$ & 总不确定度 & $v_q+\sigma_\epsilon^2$ \\
$\alpha_q$ & 直方图融合系数 & 待求 \\
$w_q\triangleq\alpha_q\hat N_q$ & 有效位置权重 & \S4 引出 \\
$b^*$ & 融合直方图的峰位 & 最终深度的来源 \\
\bottomrule
\end{tabular}
\end{center}

以下在\textbf{计数域}推导（每帧计数 $H_q$）。率域直方图 $F=H/M$ 只差公共因子 $1/M$，被系数归一化吸收，不影响 $\alpha$ 的形式。幅度一律省略因子 $M$。

\section{邻居直方图的模型：两类误差}

对齐之后，邻居 $q$ 的直方图在理想情况下是以 $\mu_q$ 为中心、幅度 $\hat N_q$ 的鼓包：
\begin{equation}
  \widetilde F_q(b)\;=\;\frac{\hat N_q}{M}\,h(b-\mu_q)\;+\;\text{底座}.
\end{equation}
实际有两类误差，性质完全不同：

\textbf{(a) 对齐误差 $\epsilon_q$——"位置错了"。}鼓包形状不变，整体中心从 $\mu_p$ 偏到 $\mu_q=\mu_p+\epsilon_q$。它是随机变量（每对邻居的对齐残差不同），方差 $\sigma_\epsilon^2$。对融合而言它是\textbf{偏置}：把鼓包搬错了地方。它作用在鼓包的\textbf{中心}上。

\textbf{(b) 计数噪声——"高度在抖"。}每个 bin 的计数是泊松型的：$\Var[H_q(b)]=\E[H_q(b)]=\lambda_q(b)$，逐 bin 独立。它让鼓包\textbf{逐 bin 高低不平}，但不移动其中心。它作用在鼓包的\textbf{高度}上。

\textbf{关键区别}：对齐误差沿 $x$ 方向搬鼓包，计数噪声沿 $y$ 方向抖高度。它们传到融合峰位 $b^*$ 上的机制不同，必须分开算（\S5）。

泊松比例（方差 = 期望，且逐 bin 独立）在率域同样成立：正文 2.2 节已证明稀疏工作区内 Coates 校正的逐 bin 方差恰好填回 $\lambda_p(b)/M^2$ 的泊松位置——所以无论在计数域还是校正后的率域推导，逐 bin 噪声与幅度成比例这一前提一致。

\section{第一步（桥）：融合峰位是鼓包中心的幅度加权平均}

融合直方图（省略 $1/M$）：
\begin{equation}
  F(b)\;=\;\sum_q \alpha_q\hat N_q\,h(b-\mu_q).
\end{equation}
定义第 $q$ 个鼓包在融合直方图中的\textbf{高度}与其\textbf{有效位置权重}：
\begin{equation}
  W_q\;\triangleq\;\alpha_q\hat N_q,
\end{equation}
则 $F(b)=\sum_q W_q\,h(b-\mu_q)$——融合直方图就是一堆\textbf{高度为 $W_q$、中心为 $\mu_q$} 的同形鼓包之和。

\subsection{一阶展开}

取任意参考点 $\bar\mu$，记 $\delta_q=\mu_q-\bar\mu$（$|\delta_q|\ll\sigma_h$）、$x=b-\bar\mu$。Taylor 展开：
\begin{equation}
  h(x-\delta_q)\;=\;h(x)-\delta_q\,h'(x)+O(\delta_q^2),
\end{equation}
代入并按 $h$、$h'$ 归并：
\begin{equation}
  F(b)\;=\;\underbrace{\Big(\sum_q W_q\Big)}_{A}\,h(x)\;-\;\underbrace{\Big(\sum_q W_q\delta_q\Big)}_{D}\,h'(x)\;+\;O(\delta^2).
\end{equation}
峰位条件 $F'(b)=0$：
\begin{equation}
  A\,h'(x)-D\,h''(x)\;=\;0.
\end{equation}
高斯 $h(x)=e^{-x^2/2\sigma_h^2}$ 满足 $h'=-\dfrac{x}{\sigma_h^2}h$、$h''=\Big(\dfrac{x^2}{\sigma_h^4}-\dfrac{1}{\sigma_h^2}\Big)h$，代入：
\begin{equation}
  -\frac{A}{\sigma_h^2}x\;-\;D\Big(\frac{x^2}{\sigma_h^4}-\frac{1}{\sigma_h^2}\Big)\;=\;0
  \quad\Longrightarrow\quad
  D\;=\;A\,x+\frac{D\,x^2}{\sigma_h^2}
  \quad\Longrightarrow\quad
  x\;=\;\frac{D}{A}\quad\big(\text{舍去 }O(\delta^2)\text{ 项}\big).
\end{equation}
于是融合峰位
\begin{equation}
  b^*\;=\;\bar\mu+\frac{D}{A}\;=\;\frac{\sum_q W_q\,\mu_q}{\sum_q W_q}
  \;=\;\boxed{\;\frac{\sum_q \alpha_q\hat N_q\,\mu_q}{\sum_q \alpha_q\hat N_q}\;}
  \label{eq:bridge}
\end{equation}

\textbf{这就是桥}：融合峰位是各邻居中心的加权平均，权重是鼓包高度 $W_q=\alpha_q\hat N_q$。\textbf{系数 $\alpha$ 自己不是位置权重——它乘上幅度之后才是。}近似条件是 $(\delta/\sigma_h)^2\ll1$：以标定值 $\sigma_\epsilon\approx0.7$、$\sigma_h\approx5$ 计，$(\delta/\sigma_h)^2\approx2\%$，充分满足。

\subsection{数字核对}

两邻居：$\hat N_1=3.5$（中心 $200.0$），$\hat N_2=35$（中心 $201.0$）。
\begin{itemize}
  \item 等系数：票数 $3.5:35=1:10$，$b^*\approx200.9$——亮邻居十倍话语权；
  \item 要同票：$\alpha_2/\alpha_1=\hat N_1/\hat N_2=1/10$——\textbf{系数必须反比于幅度}。
\end{itemize}

\section{第二步：融合峰位的方差——两条独立的噪声路}

\subsection{对齐路（$\epsilon_q$ 搬鼓包）}

式~\eqref{eq:bridge} 对任意中心位置成立。把 $\mu_q=\mu_p+\epsilon_q$ 代入：
\begin{equation}
  b^*-\mu_p\;=\;\frac{\sum_q W_q\,\epsilon_q}{\sum_q W_q}.
\end{equation}
$\epsilon_q$ 独立、方差 $\sigma_\epsilon^2$：
\begin{equation}
  \Var_{\text{align}}\;=\;\sigma_\epsilon^2\,\frac{\sum_q W_q^2}{\big(\sum_q W_q\big)^2}.
  \label{eq:var_align}
\end{equation}

\subsection{计数路（逐 bin 噪声抖高度）}

逐 bin 噪声让 $F(b)$ 上下抖动，峰位随之抖动。用位置参数的 Fisher 信息衡量"融合直方图还剩多少位置信息"：
\begin{equation}
  \mathcal I(\alpha)\;=\;\sum_b\frac{\big(\partial F(b)/\partial\mu\big)^2}{\Var[F(b)]}.
\end{equation}
两个成分逐 bin 算：
\begin{align}
  \frac{\partial F(b)}{\partial\mu} &= -\sum_q \alpha_q\hat N_q\,h'(b-\mu),\\
  \Var[F(b)] &= \sum_q \alpha_q^2\,\Var[H_q(b)]\;=\;\sum_q \alpha_q^2\,\hat N_q\,h(b-\mu),
\end{align}
末步用泊松性质（方差 $=$ 期望）；底座项的处理见 \S9 假设 2。
代入（$\beta=0$ 时分母 $\propto h(b)$，可提出求和）：
\begin{equation}
  \mathcal I(\alpha)
  =\Big(\sum_q \alpha_q\hat N_q\Big)^2
   \cdot\frac{\sum_b h'^2/h}{\sum_q \alpha_q^2\hat N_q}
  =\frac{\big(\sum_q \alpha_q\hat N_q\big)^2}{\sigma_h^2\,\sum_q \alpha_q^2\hat N_q},
\end{equation}
其中 $\sum_b h'^2/h\equiv\sum_b h'^2(b-\mu)/h(b-\mu)$。
其中用了高斯形状的恒等式 $\sum_b h'^2(x)/h(x)=1/\sigma_h^2$（与正文 2.1 节经典方差 $v=\sigma_h^2/N$ 是\textbf{同一个积分}）。于是计数路方差：
\begin{equation}
  \Var_{\text{count}}\;=\;\frac{1}{\mathcal I(\alpha)}
  \;=\;\frac{\sum_q W_q^2\,v_q}{\big(\sum_q W_q\big)^2},
  \qquad v_q=\frac{\sigma_h^2}{\hat N_q}.
  \label{eq:var_count}
\end{equation}
\textbf{核对}：等系数（$W_q\propto\hat N_q$）时 $\Var_{\text{count}}=\sigma_h^2\sum\hat N/(\sum\hat N)^2=\sigma_h^2/\sum\hat N$——正是"所有邻居的光子池化为一个鼓包"的方差 ✓。

\subsection{合计与等价说法}

两条路独立（对齐残差搬中心、计数噪声抖高度，互不相关），方差相加：
\begin{equation}
  \Var[b^*]\;=\;\frac{\sum_q W_q^2\,v^{\mathrm{eff}}_q}{\big(\sum_q W_q\big)^2},
  \qquad
  v^{\mathrm{eff}}_q\;=\;v_q+\sigma_\epsilon^2\;=\;\frac{\sigma_h^2}{\hat N_q}+\sigma_\epsilon^2.
  \label{eq:var_total}
\end{equation}
\textbf{等价说法}：式~\eqref{eq:var_total} 说明"融合直方图取峰"在方差意义下完全等价于"每个邻居先用直方图估计一次中心（方差 $v^{\mathrm{eff}}_q$），再以 $W_q=\alpha_q\hat N_q$ 为权做加权平均"。两种架构描述同一件事——这就是"位置坐标"与"直方图坐标"的关系。

\section{第三步：方差极小化}

式~\eqref{eq:var_total} 对 $W$ 整体缩放不变（分子分母都是二次齐次），固定 $\sum W_q=1$，极小化 $\sum_q W_q^2v^{\mathrm{eff}}_q$。拉格朗日乘子：$2W_qv^{\mathrm{eff}}_q=\lambda$，即
\begin{equation}
  W_q\;\propto\;\frac{1}{v^{\mathrm{eff}}_q}
  \;=\;\frac{1}{\sigma_h^2/\hat N_q+\sigma_\epsilon^2}
  \;=\;\frac{\hat N_q}{\sigma_h^2+\sigma_\epsilon^2\hat N_q}.
  \label{eq:W_opt}
\end{equation}
教科书结论：逆方差加权。注意 $W$ 的形态——随 $\hat N_q$ \textbf{递增、饱和于 $1/\sigma_\epsilon^2$}。两个极限：
\begin{itemize}
  \item $\sigma_\epsilon\to0$（无对齐残差）：$W\propto\hat N_q/\sigma_h^2$，仍递增；回代 $\alpha_q=W_q/\hat N_q=$ 常数——\textbf{等系数}。等权池化正是纯计数情形的联合最优（数值验证：Fisher 比值 $1.000$）。
  \item $\hat N_q\gg\hat N_0$：$W_q\to1/\sigma_\epsilon^2$——\textbf{封顶}。多余的光子不再改善这一邻居的位置知识，因为对齐残差是瓶颈。
\end{itemize}

\section{第四步：回到直方图系数}

由 $W_q=\alpha_q\hat N_q$：
\begin{equation}
  \alpha_q\;=\;\frac{W_q}{\hat N_q}
  \;\propto\;\frac{1}{\sigma_h^2+\sigma_\epsilon^2\hat N_q}
  \;=\;\frac{1}{\sigma_\epsilon^2}\cdot\frac{1}{\hat N_q+\sigma_h^2/\sigma_\epsilon^2}.
\end{equation}
常数 $1/\sigma_\epsilon^2$ 被系数归一化吸收：
\begin{equation}
  \boxed{\;\alpha_q\;\propto\;\frac{1}{\hat N_q+\hat N_0},
  \qquad \hat N_0\;\triangleq\;\frac{\sigma_h^2}{\sigma_\epsilon^2}.\;}
\end{equation}
\textbf{$\hat N_0$ 的物理含义}：使计数方差 $\sigma_h^2/\hat N$ 与对齐方差 $\sigma_\epsilon^2$ 相等的光子数——"计数受限"与"对齐受限"的交叉点。
\begin{itemize}
  \item $\hat N_q\ll\hat N_0$：计数受限。中心精度随光子数改善，系数接近常数（等权池化近优）；
  \item $\hat N_q\gg\hat N_0$：对齐受限。位置精度被 $\sigma_\epsilon$ 封顶，多余幅度不携带位置信息，系数必须反比下降——否则亮邻居的"虚精度"支配融合。
\end{itemize}
以正文标定值 $\hat N_0=50$、$\sigma_h\approx5.03$ bin（TDC 口径物理值）换算：$\sigma_\epsilon=\sigma_h/\sqrt{\hat N_0}\approx0.71$ bin——亚 bin 对齐残差，物理合理。又，典型峰显著度 $\hat N_q\sim3.5\text{--}10\ll\hat N_0$：多数邻居处于计数受限区，系数近似常数、轻度均衡——这解释了工程上该因子作用温和。

\textbf{常见误解澄清}："幅度因子替代精度因子"的说法是误读。精度因子是位置权重 $W$（递增饱和）；$\alpha$ 是 $W$ 除以幅度后的直方图坐标形态（递减）。两者由 $W=\alpha\hat N$ 定义性地联系——推导中 $1/v^{\mathrm{eff}}\propto\hat N$（递增）自始至终未被违反；$\alpha$ 里的 $1/\hat N$ 是换算除法，不是对精度方向的断言。

\section{数值验证}

\textbf{设置}：两邻居 $\hat N_1=3.5$、$\hat N_2=70$；$\sigma_h=5$ bin；$\sigma_\epsilon=0.7071$ bin（$\hat N_0=50$）；$\beta=0$。每次试验：$\epsilon_q\sim\mathcal N(0,\sigma_\epsilon^2)$ 独立重采样 $\to$ 泊松抽样 $\to$ 融合 $F=H_1+rH_2$（$\alpha_2/\alpha_1=r$）$\to$ 匹配滤波 + argmax + 三点抛物线亚 bin。每点 $2\times10^5$ 次试验。

解析曲线（式~\eqref{eq:var_total}）：$\Var(r)=\dfrac{W_1^2v^{\mathrm{eff}}_1+W_2^2v^{\mathrm{eff}}_2}{(W_1+W_2)^2}$，$W_1=3.5$、$W_2=70r$。预测极小 $r^*=\dfrac{\hat N_1+\hat N_0}{\hat N_2+\hat N_0}=0.446$。

\begin{center}
\begin{tabular}{@{}rrr@{}}
\toprule
$r=\alpha_2/\alpha_1$ & MC 方差 & 解析方差 \\
\midrule
0.150 & 1.3312 & 0.9598 \\
0.250 & 1.0698 & 0.8075 \\
0.350 & 0.9958 & 0.7757 \\
0.446 & 0.9802 & 0.7707 \\
0.550 & \textbf{0.9767} & 0.7733 \\
0.700 & 0.9771 & 0.7806 \\
0.900 & 0.9846 & 0.7905 \\
1.200 & 0.9996 & 0.8022 \\
1.600 & 1.0053 & 0.8130 \\
2.200 & 1.0250 & 0.8232 \\
\bottomrule
\end{tabular}
\end{center}

\textbf{结果}：(1) MC 方差曲线与解析曲线\textbf{形状一致}，且都在 $r\approx0.45$--$0.55$ 处取极小——MC 中 $r=0.446$ 与 $0.55$ 差异仅 $1.1\sigma$（统计不可区分），而 $r=0.25$ 差 $+9.5\%$、$r=0.9$ 差 $+0.8\%$（$2.5\sigma$）、$r\ge1.2$ 差 $2\%$ 以上，均显著更差；预测 $r^*=0.446$ 落在 MC 极小平坦区内。(2) MC 绝对方差高于解析曲线约 $24\%$--$39\%$：解析曲线的计数路以 CRLB 为准，稀疏计数（$\hat N_1=3.5$）下取峰估计未完全有效；缺口随 $r$ 平缓变化，不影响极小位置。(3) 极小平坦本身是物理结论：邻域内邻居的信背比相近，权重形状的选择本来就不敏感——这也是正文标定温和的原因。

\textbf{极限一致性}：$\sigma_\epsilon=0$ 时式~\eqref{eq:W_opt} 给出 $W\propto\hat N_q$、$\alpha_q=$ 常数（等权）；纯计数 Fisher 分析（Rayleigh 商，精确）与数值验证均给出等权最优、且达到联合 ML 上界（比值 $1.000$）——与本文档框架的 $\hat N_0\to\infty$ 极限完全一致。

\section{与正文三章的对应}

\begin{center}
\small
\begin{tabular}{@{}lll@{}}
\toprule
本文档 & 正文位置 & 承担内容 \\
\midrule
\S3 两类误差 & 2.4（$\epsilon$ 分解）+ 2.2（计数噪声） & 偏置 vs 方差的分工 \\
\S4 桥（式~\eqref{eq:bridge}） & 2.4 需补的直方图系数推导 & 系数 $\leftrightarrow$ 位置权重换算 \\
\S4 $v^{\mathrm{eff}}$、\S5 $W\propto1/v^{\mathrm{eff}}$ & 2.4（总不确定度、逆方差权重） & 理想权重 \\
$\gamma$ 乘入 $W\propto\gamma/v^{\mathrm{eff}}$ & 2.5（可融合性概率） & 跨表面粗差甄别 \\
\S6 $\alpha\propto1/(\hat N+\hat N_0)$ & 2.6（工程实现） & $\gamma\to g_\Delta g_\rho$，$\hat N\to$ 峰显著度 \\
\bottomrule
\end{tabular}
\end{center}

分工要点：$\sigma_\epsilon$（经 $\hat N_0$）只刻画\textbf{同一表面内}的对齐残差尺度，是可标定的常数；"这一对邻居是否根本不在同一表面"是粗差，$\sigma_\epsilon$ 刻画不了——由 $\gamma$（几何 + 互相关证据）保守处理。

\section{假设清单}

\begin{enumerate}
  \item 高斯 IRF。一般形状下 $\sigma_h^2\to\big(\sum_bh'^2/h\big)^{-1}$，结构性结论不变；
  \item 稀疏工作区：Coates 校正无损，逐 bin 方差 $\propto$ 幅度（正文 2.2）；峰窗内背景项可略（$\beta>0$ 时 Fisher 分母不再 $\propto h$，清洁形式退化，为高阶修正）；
  \item 小位移：$\sigma_\epsilon\ll\sigma_h$（标定值下 $(\sigma_\epsilon/\sigma_h)^2\approx2\%$）；
  \item 取峰近似有效（CRLB 可达）。稀疏计数下实测有 $\sim25\%$--$40\%$ 效率缺口（\S7），缺口随 $r$ 平缓，不改变极小位置；
  \item 邻居间独立（对齐残差、计数噪声均独立）。
\end{enumerate}

\end{document}
